\documentclass{article}
\usepackage[utf8]{inputenc}
\usepackage[russian]{babel}
\usepackage{amsmath}

\begin{document}

\title{Домашнее задание №2 по обработке сигналов}
\author{}
\date{}
\maketitle

\section{Задача 1}

Порождающая матрица расширенного кода Хэмминга с параметрами $(8,4)$ в систематической форме имеет следующий вид:

\[
G = \begin{pmatrix}
1 & 0 & 0 & 0 & 0 & 1 & 1 & 1 \\
0 & 1 & 0 & 0 & 1 & 0 & 1 & 1 \\
0 & 0 & 1 & 0 & 1 & 1 & 0 & 1 \\
0 & 0 & 0 & 1 & 1 & 1 & 1 & 0
\end{pmatrix}
\]

Для этого кода минимальное расстояние составляет $d_{\min} = 4$.

Анализ весов всех кодовых слов показывает: существует ровно одно слово нулевого веса, 14 слов с весом 4 и одно слово с весом 8 (слово из одних единиц).

Благодаря линейности кода расстояния между различными кодовыми словами совпадают с весами их разностей, а распределение весов разностей, очевидно, такое же, как распределение весов самих слов: для пары одинаковых слов расстояние равно 0, для всех остальных пар — либо 4, либо 8.

Для проверки распределения весов был использован следующий скрипт на языке Python:

\begin{verbatim}
import numpy as np

G = np.array([
[1,0,0,0,0,1,1,1],
[0,1,0,0,1,0,1,1],
[0,0,1,0,1,1,0,1],
[0,0,0,1,1,1,1,0]
], dtype=int) % 2

weights = []
for i in range(16):
    u = np.array([int(b) for b in format(i, '04b')])
    c = (u @ G) % 2
    wt = np.sum(c)
    weights.append(wt)

unique_weights = set(weights)
min_d = min(w for w in weights if w > 0)
print(f'Unique weights: {unique_weights}')
print(f'Min distance: {min_d}')
\end{verbatim}

Выполнение кода даёт следующий результат: Unique weights: $\{0, 8, 4\}$, Min distance: 4. Таким образом, получаем одно слово веса 0, 14 слов веса 4 и одно слово веса 8.

\section{Задача 2}

Проверочная матрица кода Хэмминга $(7,4)$ задаётся как

\[
H = \begin{pmatrix}
0 & 0 & 0 & 1 & 1 & 1 & 1 \\
0 & 1 & 1 & 0 & 0 & 1 & 1 \\
1 & 0 & 1 & 0 & 1 & 0 & 1
\end{pmatrix}
\]

Порождающая матрица дуального кода получается непосредственно из $H$ (с точностью до перестановки столбцов, которая не изменяет весовые характеристики). Следовательно, дуальный код имеет параметры $(7,3)$.

Такой код представляет собой симплекс-код длины 7, который может быть построен из всех ненулевых двоичных векторов длины 3, расположенных в качестве столбцов порождающей матрицы. Для симплекс-кода $(7,3)$ известно, что все ненулевые кодовые слова обладают весом $2^{3-1} = 4$.

Убедимся в этом с помощью аналогичного расчёта на Python:

\begin{verbatim}
import numpy as np

G_star = np.array([
[0,0,0,1,1,1,1],
[0,1,1,0,0,1,1],
[1,0,1,0,1,0,1]
], dtype=int) % 2

weights = []
for i in range(8):
    u = np.array([int(b) for b in format(i, '03b')])
    c = (u @ G_star) % 2
    wt = np.sum(c)
    weights.append(wt)

print(set(weights))
\end{verbatim}

Результат выполнения: $\{0, 4\}$, что подтверждает, что каждое ненулевое кодовое слово дуального кода имеет вес 4, т.е. код действительно является симплекс-кодом.

\section{Задача 3}

Рассмотрим код с проверкой на чётность длины $n$. Его параметры: $(n, n-1)$, проверочная матрица представляет собой строку из всех единиц: $H = (1, 1, \dots, 1)$.

Порождающая матрица дуального кода совпадает с $H$, но поскольку ранг $H$ равен единице, дуальный код является $(n, 1)$-кодом повторения: он содержит только два кодовых слова — нулевое и слово из всех единиц.

Таким образом, дуальный код имеет следующие характеристики:
\begin{itemize}
    \item длина $n$ (та же, что и у исходного кода);
    \item размерность $k = 1$;
    \item число кодовых слов $2^k = 2$;
    \item минимальное расстояние $d_{\min} = n$ (расстояние между единственными ненулевыми словами — нулевым и словом из всех единиц).
\end{itemize}

Данный код способен исправлять любые комбинации ошибок кратности не более $\lfloor (n-1)/2 \rfloor$.

\end{document}